\documentclass[11pt]{article}
\usepackage[a4paper,margin=28mm]{geometry}
\usepackage{amsmath,amssymb,amsthm}
\usepackage[T1]{fontenc}
\usepackage{lmodern}
\usepackage{microtype}
\usepackage[colorlinks=true,linkcolor=blue,citecolor=blue,urlcolor=blue]{hyperref}
\hypersetup{pdftitle={A counterexample to degree-zero universality of the interleaving distance},pdfauthor={Kaiwen Jin}}
\newtheorem{theorem}{Theorem}
\newtheorem{proposition}[theorem]{Proposition}
\newtheorem{lemma}[theorem]{Lemma}
\newcommand{\R}{\mathbb R}
\newcommand{\one}{\mathbf 1}
\newcommand{\colim}{\mathop{\mathrm{colim}}}
\newcommand{\im}{\mathop{\mathrm{im}}}
\newcommand{\supp}{\mathop{\mathrm{supp}}}
\newcommand{\Vect}{\mathrm{Vect}}
\title{A counterexample to degree-zero universality\texorpdfstring{\\}{ }of the interleaving distance}
\author{Kaiwen Jin\thanks{\href{mailto:yui.ui.siki@gmail.com}{yui.ui.siki@gmail.com}}}
\date{13 October 2026}
\begin{document}
\maketitle

\begin{abstract}
The universality conjecture for the interleaving distance asks whether every
distance stable under uniform perturbations of sublevel functions is bounded
by the interleaving distance on realizable persistence modules. Its formulation
for each homology degree is false for ordinary zero-dimensional homology.
For every field we construct a bounded, $0$-stable pseudometric on all
four-parameter persistence modules and a pair realized on three-point discrete
spaces for which this distance is $3$ and the interleaving distance is $1$.
The algebraic pair is Bjerkevik's published free-module example. The additional
argument recovers the upper supports of eventual sublevel components from the
image of a module in its colimit and proves stability for arbitrary spaces and
functions, with no finite-dimensionality assumption.
\end{abstract}

\section{The degree-zero question}

Sublevel persistent homology records how the topology of a space changes as
one or several thresholds increase. A distance between the resulting modules
is stable when a uniform change in the threshold function bounds the change
in the descriptor. Universality asks whether the interleaving distance is the
largest distance with this property on modules that arise from homology.
Lesnick proved this assertion over prime fields in positive homology degrees
\cite[Theorem 5.5]{Lesnick}. His Conjecture 5.7 includes degree zero and
arbitrary fields, with the comparison restricted to realizable modules.
We give a counterexample to its ordinary degree-zero clause.

Fix a field $k$. All modules in this note are functors
$M:(\R^4,\leq)\to\Vect_k$, with coordinatewise order. Write
$\one=(1,1,1,1)$ and $M(\varepsilon)_a=M_{a+\varepsilon\one}$.
An $\varepsilon$-interleaving, $\varepsilon\geq0$, consists of natural maps
$\varphi:M\to N(\varepsilon)$ and $\psi:N\to M(\varepsilon)$ satisfying
\[
 \psi(\varepsilon)\varphi=\eta_M^{2\varepsilon},\qquad
 \varphi(\varepsilon)\psi=\eta_N^{2\varepsilon},
\]
where $\eta_M^{2\varepsilon}$ is the structure map
$M_a\to M_{a+2\varepsilon\one}$. The distance $d_I$ is the infimum of
such $\varepsilon$, or $+\infty$ if none exists.
For any topological space $X$ and arbitrary function $f:X\to\R^4$, let
\[
 P_0(f)_a=H_0\bigl(\{x\in X:f(x)\leq a\};k\bigr),
\]
using ordinary singular homology and sublevel inclusions. A pseudometric on
module isomorphism classes is $0$-stable if
\begin{equation}\label{eq:stability}
 d(P_0(f),P_0(g))\leq
 \sup_{x\in X}\|f(x)-g(x)\|_\infty
\end{equation}
for every $X,f,g$, with the empty supremum equal to zero. Extended
pseudometrics may take the value $+\infty$ and may vanish on distinct classes.

\begin{theorem}\label{thm:main}
For every field $k$, there is a $[0,3]$-valued pseudometric $d_0$ on
isomorphism classes of all four-parameter $k$-modules satisfying
\eqref{eq:stability}. There are continuous functions $f,g$ on the discrete
space $\{1,2,3\}$ such that, for $M=P_0(f)$ and $N=P_0(g)$,
\[
 d_I(M,N)=1,\qquad d_0(M,N)=3.
\]
\end{theorem}

The finite algebraic separation used here is Bjerkevik's example
\cite[Examples 5.2--5.3]{Bjerkevik}. A barcode matching distance on free modules
alone does not supply the required distance on all module classes. We construct
such a distance from the supports of components in the directed colimit.
Uniform perturbations match these components, whereas linear interleavings
can mix their basis vectors. This distinction gives the inequality in
Theorem~\ref{thm:main}.

The conclusion concerns ordinary $H_0$. It does not resolve the
arbitrary-field question in positive degrees, nor does it contradict
prime-field universality for a distance required to be stable in every
homology degree \cite[Corollary 5.6]{Lesnick}. Lesnick's one-parameter
pointwise finite-dimensional result in Section 5.4 uses reduced homology;
that convention is not used here.

\section{Supports in the image of a module's colimit}

For any module $M$, put
\[
 C_M=\colim_{a\in\R^4}M_a,\qquad
 I_M(a)=\im(M_a\longrightarrow C_M).
\]
The subspaces $I_M(a)$ increase with $a$, and their union is $C_M$.
Indeed, the vector-space colimit is the direct sum of the $M_a$ modulo
structure-map relations. An element has finitely many representatives, which
can be transported to a common upper grade.

Consider modules for which $C_M$ is finite-dimensional and has a basis
$e_1,\ldots,e_r$ such that every $I_M(a)$ is spanned by a subset of this
basis. The zero colimit is included. For each basis vector define its upper
support
\[
 U_j=\{a\in\R^4:e_j\in I_M(a)\}.
\]
An upper set is a subset closed under increasing coordinates. Each $U_j$
is nonempty because the images exhaust $C_M$. Let $\mathcal S(M)$ be the
finite multiset of these supports, with one occurrence for each basis vector.

\begin{proposition}\label{prop:intrinsic}
The existence of such a coordinate basis and the multiset
$\mathcal S(M)$ depend only on the module's isomorphism class.
\end{proposition}
\begin{proof}
We recover each support multiplicity directly from the subspaces $I_M(a)$.
For a nonempty upper set $U$, define
\[
 K_U=\bigcap_{a\in U}I_M(a),\qquad
 L_U=\sum_{\substack{V\text{ nonempty upper}\ U\subsetneq V}}K_V.
\]
The second subspace lies in the first. In a coordinate basis, membership
in any intersection can be tested coefficient by coefficient, giving
\[
 K_U=\operatorname{span}\{e_j:U\subseteq U_j\},\qquad
 L_U=\operatorname{span}\{e_j:U\subsetneq U_j\}.
\]
For the second identity, a basis vector belonging to $K_V$ has support
containing $V$ and therefore properly containing $U$. Conversely, if
$U\subsetneq U_j$, that vector belongs to the summand $K_{U_j}$ of $L_U$.
Thus
\[
 \dim(K_U/L_U)=\#\{j:U_j=U\}.
\]
This recovers the multiplicity of every upper set without choosing a basis.
A natural module isomorphism induces a colimit isomorphism carrying every
$I_M(a)$ to the corresponding image. It therefore preserves the coordinate
basis property and all the recovered multiplicities.
\end{proof}

\section{A distance stable for all ordinary degree-zero inputs}

For nonempty upper sets $U,V$, define
\[
 \rho(U,V)=\inf\{\varepsilon\geq0:
 U+\varepsilon\one\subseteq V,\quad
 V+\varepsilon\one\subseteq U\}.
\]
The infimum of the empty set is $+\infty$. This is an extended pseudometric.
To check the triangle inequality, two pairs of shifted containments compose
with the shifts added. Taking arbitrarily close upper bounds for the two
infima gives the triangle inequality without an attainment assumption.

For finite multisets $\mathcal U,\mathcal V$ with the same number of
occurrences, let
\[
 \beta(\mathcal U,\mathcal V)
   =\inf_{\sigma}\max_{U\in\mathcal U}\rho(U,\sigma(U)),
\]
where $\sigma$ ranges over occurrence bijections. Set the distance to
$+\infty$ for unequal cardinalities and set
$\beta(\varnothing,\varnothing)=0$. Composing bijections and applying the
triangle inequality for $\rho$ proves that $\beta$ is an extended
pseudometric. Adjoin a symbol $*$ at distance $+\infty$ from every multiset
and at distance zero from itself; write $\beta_*$ for this extension.

Define a feature on every module class by
\[
 \mathcal F(M)=
 \begin{cases}
 \mathcal S(M),& C_M\text{ has a finite coordinate basis as above},\\
 *,&\text{otherwise},
 \end{cases}
\]
Then put
\[
 d_0(M,N)=\min\{3,\beta_*(\mathcal F(M),\mathcal F(N))\}.
\]
Proposition~\ref{prop:intrinsic} makes this definition independent of all
choices. Pullback preserves the pseudometric properties, and the cap preserves
the triangle inequality by
$\min\{3,s+t\}\leq\min\{3,s\}+\min\{3,t\}$ for
$s,t\in[0,+\infty]$. Consequently $d_0$ is a bounded pseudometric on the
entire class of modules.

\begin{proposition}\label{prop:stable}
The distance $d_0$ is $0$-stable for every topological space and arbitrary
functions to $\R^4$.
\end{proposition}
\begin{proof}
Ordinary singular $H_0(Y;k)$ is naturally the free vector space
$k[\pi_0(Y)]$ on path components: singular $0$-chains are free on points,
and boundaries of singular $1$-simplices impose precisely the relations
identifying endpoints of paths. This description is natural for maps.
Write
\[
 \Pi_a^f=\pi_0\bigl(\{x:f(x)\leq a\}\bigr),\qquad
 B_f=\colim_a\Pi_a^f
\]
with the second colimit taken in sets. The free-vector-space construction
commutes with this colimit. For any vector space $W$, compatible linear maps
from $k[\Pi_a^f]$ to $W$ are exactly compatible set maps from $\Pi_a^f$ to
the underlying set of $W$, hence set maps from $B_f$ to $W$. This universal
property gives
\[
 C_{P_0(f)}\cong k[B_f].
\]
The image $I_{P_0(f)}(a)$ is spanned by the subset of $B_f$ reached from
$\Pi_a^f$. If $B_f$ is finite, it supplies a coordinate basis and its supports
are exactly $\mathcal S(P_0(f))$. If $B_f$ is infinite, the colimit is
infinite-dimensional and $\mathcal F(P_0(f))=*$. We use this actual sublevel
colimit throughout; no identification of $B_f$ with $\pi_0(X)$ is needed.

Let $\delta=\sup_x\|f(x)-g(x)\|_\infty$. If $\delta=+\infty$, the required
bound is automatic. Otherwise the sublevel inclusions induce natural set maps
\[
 \Pi_a^f\longrightarrow\Pi_{a+\delta\one}^g,
 \qquad
 \Pi_a^g\longrightarrow\Pi_{a+\delta\one}^f.
\]
Their composites are the structure maps at shift $2\delta$. They therefore
induce inverse bijections $\alpha:B_f\to B_g$ on the colimits: the composite
sends any represented component to its later structure image, representing
the same colimit element. Compatibility of the maps makes this independent
of the chosen representative.

If these sets are infinite, both features are $*$ and $d_0=0$. If finite,
$\alpha$ matches their supports with
\[
 U_b^f+\delta\one\subseteq U_{\alpha(b)}^g,
 \qquad
 U_{\alpha(b)}^g+\delta\one\subseteq U_b^f.
\]
The first containment follows from the first component map at every grade
in $U_b^f$. The reverse map gives the second, since its colimit bijection
is $\alpha^{-1}$. Thus each matched $\rho$ is at most $\delta$, so
$\beta\leq\delta$ and $d_0\leq\delta$. The argument includes $\delta=0$.
For empty $X$ both multisets are empty and $d_0=0$.
\end{proof}

In particular, finite perturbations preserve the finite or infinite
cardinality of the component colimit. This is why assigning the exceptional
feature to infinite-colimit inputs is compatible with the stability bound.
Individual sublevel homology spaces may still be infinite-dimensional when
the component colimit is finite.

\section{The three-point counterexample}

For $p\in\R^4$, let $k_{\uparrow p}$ denote the module with value $k$ at
$a\geq p$, zero elsewhere, and identity maps between nonzero values.
Take the birth vectors of Bjerkevik's four-parameter free-module example
\cite[Example 5.3]{Bjerkevik}:
\[
\begin{aligned}
 u_1&=(0,1,10,11),&v_1&=(1,0,11,10),\\
 u_2&=(0,-1,12,11),&v_2&=(1,0,9,12),\\
 u_3&=(2,1,10,9),&v_3&=(-1,2,11,10).
\end{aligned}
\]
Put $M=\bigoplus_{i=1}^3k_{\uparrow u_i}$ and
$N=\bigoplus_{j=1}^3k_{\uparrow v_j}$. On the discrete space
$X=\{1,2,3\}$, define $f(i)=u_i$ and $g(i)=v_i$. Its sublevels consist of
the points whose birth vectors lie below the grade, so naturally
$P_0(f)\cong M$ and $P_0(g)\cong N$. Both functions are continuous.

A morphism $k_{\uparrow p}\to k_{\uparrow q}(\varepsilon)$ is a scalar map
if $p+\varepsilon\one\geq q$, and otherwise it is zero. Indeed, in the first
case every source grade lies in the shifted target, and naturality at common
upper grades forces one scalar. In the second case the map at $p$ is zero;
naturality from $p$ and any other source grade to a common upper grade forces
every component to be zero. This also describes each component of a morphism
between the finite direct sums.

The least directed shifts for $M_i\to N_j(\varepsilon)$ and
$N_j\to M_i(\varepsilon)$ are respectively
\[
 T^{\to}=\begin{pmatrix}1&1&1\\1&1&3\\1&3&1\end{pmatrix},\qquad
 T^{\leftarrow}=\begin{pmatrix}1&1&1\\1&3&1\\1&1&3\end{pmatrix}.
\]
The first matrix has source $M_i$ in row $i$ and target $N_j$ in column $j$;
the second has source $N_j$ in row $j$ and target $M_i$ in column $i$.
Their entries are obtained by taking the maximum coordinate of the
target birth vector minus the source birth vector.

The inverse matrices from Bjerkevik's Example 5.2, equation (15), are
\[
 A=\begin{pmatrix}1&1&1\\1&1&0\\1&0&1\end{pmatrix},\qquad
 B=\begin{pmatrix}-1&1&1\\1&0&-1\\1&-1&0\end{pmatrix}.
\]
Use the source-row convention: the component of $\varphi:M\to N(1)$ from
$M_i$ to $N_j(1)$ has coefficient $A_{ij}$, and the component of
$\psi:N\to M(1)$ from $N_j$ to $M_i(1)$ has coefficient $B_{ji}$.
Every nonzero entry has directed threshold
$1$, so these are natural transformations. Direct integer multiplication gives
\[
 AB=
 \begin{pmatrix}
 -1+1+1&1-1&1-1\\
 -1+1&1&1-1\\
 -1+1&1-1&1
 \end{pmatrix}=I_3,\qquad BA=I_3.
\]
Each nonzero composite component follows two allowed support containments.
Thus at every active source grade the composite is given by the same constant
matrix product, and these products give the structure maps at shift $2$.
The integer identities hold over every field, including characteristic $2$.
We have a $1$-interleaving.

For $0\leq\varepsilon<1$, every component map $M_i\to N_j(\varepsilon)$
is zero, as all entries of $T^{\to}$ are at least $1$. The map
$\varphi$ must therefore be zero. It cannot compose to
$\eta_M^{2\varepsilon}$, which is nonzero at any grade with nonzero $M$.
Consequently $d_I(M,N)=1$.

The image families of $M$ and $N$ in their three-dimensional colimits are
coordinate families. Their support multisets consist of the principal upper
sets $\uparrow u_i$ and $\uparrow v_j$. The shifted-containment definition gives
\[
 \rho(\uparrow p,\uparrow q)=\|p-q\|_\infty.
\]
The matching-cost matrix is therefore
\[
 \bigl(\|u_i-v_j\|_\infty\bigr)_{i,j}
 =\begin{pmatrix}1&1&1\\1&3&3\\1&3&3\end{pmatrix}.
\]
A matching of cost below $3$ would have to send both the second and third
$M$ occurrences to the first $N$ occurrence. This is impossible for a
bijection. Every entry is at most $3$, so a matching of cost $3$ exists.
It follows that $d_0(M,N)=3$. Together with Proposition~\ref{prop:stable},
this proves Theorem~\ref{thm:main}: the realizable ordinary-$H_0$ modules above
have stable distance $3$ and interleaving distance $1$, over every field.

\begin{thebibliography}{9}
\bibitem{Lesnick}
M. Lesnick, \emph{The Theory of the Interleaving Distance on Multidimensional
Persistence Modules}, Foundations of Computational Mathematics
\textbf{15} (2015), 613--650.
\href{https://doi.org/10.1007/s10208-015-9255-y}{doi:10.1007/s10208-015-9255-y}.
Author version: \href{https://arxiv.org/abs/1106.5305v4}{arXiv:1106.5305v4}.

\bibitem{Bjerkevik}
H. B. Bjerkevik, \emph{Stability of higher-dimensional interval decomposable
persistence modules}, \href{https://arxiv.org/abs/1609.02086v3}{arXiv:1609.02086v3},
21 January 2020; Examples 5.2--5.3, pp. 17--18.
\end{thebibliography}

\section*{AI-use statement}
This manuscript was prepared with assistance from OpenAI Codex for source
retrieval, proof development, drafting, and checks of the finite integer
calculations. The mathematical argument was reviewed by the same AI assistant.
Independent human mathematical review is requested.

\end{document}
